% =============================================================================
% Title:        Instrument notes
% Author:       Nathaniel Thomas
% Affiliation:  Texas A&M University, Department of Nuclear Engineering
% Email:        nathaniel@tamu.edu
% Date Created: 3/7/2026
% Version:      1.0
%
% Description:
%   A .tex source file describing the instruments used in the FLiNaK purifier
%   and some important operating details.
%
% =============================================================================

\section{Instrument Notes}

\subsection{Honeywell Satellite XT (HF Detector)}

The Honeywell Satellite XT is used to monitor the concentration of HF
inside of the fumehood of the purifier. The concentration of HF inside of the
fumehood is monitored during purifier operation.

% TODO: Implement this
% If the concentration increases above zero, the LabView software
% automatically sends
% a notification to (512)-952-1017 (Nathaniel Thomas's personal phone
% number). It
% also sends an email to nathaniel@tamu.edu.

% If the concentration of HF reaches the TLV (3 ppm from
% https://www.cdc.gov/niosh/idlh/7664393.html),
% then the LabView program calls Nathaniel Thomas at (512)-952-1017
% and ceases HF flow, pausing the
% purification. It also sends an email to nathaniel@tamu.edu.

% If the concentration of HF reaches double the TVL (6 ppm), the LabView
% program calls Nathaniel Thomas at (512)-952-1017 and shuts down all gas flow
% except argon, and shuts off all heating. It also sends an email to
% nathaniel@tamu.edu.

% If the concentration of HF reaches IDLH (30 ppm), the LabView program shuts
% down all gas flow and heating. It then calls 911 and notifies the 911 operator
% of the location, condition, and emergency. Then it notifies EHS's
% emergency line at 979-845-4311. It then calls Nathaniel Thomas at
% (512)-952-1017 and notifies of the situation. It also sends an
% email to nathaniel@tamu.edu

\subsection{SRS RGA100}

The SRS RGA100 monitors the composition of the scrubber gas. This can
be used to monitor the progress of the purification by detecting the
quantity of hydrogen,
oxygen, water, \ch{HF}, and other product gases from the purification. \par

Note that the manual does not permit operation in corrosive gas environments,
which is exactly how it is being used here. \par
% TODO: Consult SRS on this use case.

Also note, there is the possibility of cross-talk with a hot ion
gauge, which is also used very close to the detector. The manual
states that this interference can be minimized with appropriate shielding.
% TODO: Confirm that the ion gauge can (or cannot) be operated in the vicinity
% of the ion gauge.

\subsection{NI USB-6353}

% TODO

% \subsection{}

\pagebreak
